\chapter{Lesson 4: Transportation}

\section{Vocabulary}

\begin{center}
\begin{tabular}{llll}
\toprule
Characters & Pinyin & Part of speech & Meaning \\
\midrule
\zh{坐} & zuò & verb & to ride (as a passenger), take \\
\zh{开车} & kāi chē & verb-object & to drive \\
\zh{骑} & qí & verb & to ride (astride) \\
\zh{公共汽车} & gōnggòng qìchē & noun & bus \\
\zh{公交车} & gōngjiāochē & noun & bus (more common now) \\
\zh{地铁} & dìtiě & noun & subway \\
\zh{火车} & huǒchē & noun & train \\
\zh{高铁} & gāotiě & noun & high-speed rail \\
\zh{飞机} & fēijī & noun & airplane \\
\zh{出租车} & chūzūchē & noun & taxi \\
\zh{自行车} & zìxíngchē & noun & bicycle \\
\zh{走路} & zǒu lù & verb-object & to walk \\
\zh{机场} & jīchǎng & noun & airport \\
\zh{车站} & chēzhàn & noun & station, stop \\
\zh{站} & zhàn & noun / MW & stop; (number of) stops \\
\zh{线} & xiàn & noun & (subway) line \\
\zh{换车} & huàn chē & verb-object & to transfer \\
\zh{上车} / \zh{下车} & shàng chē / xià chē & verb-object & get on / get off \\
\zh{麻烦} & máfan & adjective / verb & troublesome; to trouble \\
\zh{从} & cóng & preposition & from \\
\zh{到} & dào & verb / prep. & to arrive; to \\
\zh{要是\ldots{}就\ldots} & yàoshi \ldots jiù \ldots & & if \ldots then \ldots \\
\zh{一\ldots{}就\ldots} & yī \ldots jiù \ldots & & as soon as \ldots \\
\zh{每} & měi & pronoun & every \\
\zh{堵车} & dǔ chē & verb-object & traffic jam \\
\zh{寒假} & hánjià & noun & winter vacation \\
\zh{票} & piào & noun & ticket \\
\bottomrule
\end{tabular}
\end{center}

\section{Means of transport}

\begin{important}[Which verb for which vehicle]
\begin{itemize}
  \item \zh{坐} (sit as a passenger): \zh{坐公交车, 坐地铁, 坐飞机, 坐火车, 坐出租车}.
  \item \zh{开} (operate): \zh{开车, 开飞机}.
  \item \zh{骑} (sit astride): \zh{骑自行车, 骑马, 骑摩托车}.
  \item \zh{走路} on foot; \zh{打车} to hail a taxi (or order one by app).
\end{itemize}
\end{important}

\subsection{Means-of-transport word order}

``How you go'' comes \emph{before} ``where you go''.

\begin{center}
\begin{tikzpicture}[w/.style={draw, rounded corners, minimum height=0.85cm, align=center, font=\small}]
  \node[w, fill=blue!10] (s) at (0,0) {\zh{我}\\ subject};
  \node[w, fill=purple!10] (t) at (1.9,0) {\zh{明天}\\ time};
  \node[w, fill=green!15] (m) at (4.1,0) {\zh{坐高铁}\\ means};
  \node[w, fill=orange!20] (v) at (6.3,0) {\zh{去上海}\\ \zh{去/来} + place};
  \node[w, fill=yellow!25] (p) at (8.8,0) {\zh{看朋友}\\ purpose};
  \foreach \a/\b in {s/t,t/m,m/v,v/p} \draw[->, thick] (\a) -- (\b);
  \node[font=\small] at (4.4,-1.1) {``Tomorrow I'm taking the high-speed train to Shanghai to see a friend.''};
\end{tikzpicture}
\end{center}

\begin{itemize}
  \item \zh{你每天怎么来学校？}---\zh{我骑自行车来。} How do you get to school every day?
  \item \zh{寒假你打算怎么回家？}---\zh{先坐飞机到北京，再坐火车回家。}
\end{itemize}

\section{Grammar}

\subsection{\zh{从\ldots{}到\ldots}}

\begin{itemize}
  \item \zh{从我家到学校很远。} It's far from my home to school.
  \item \zh{从滑铁卢到多伦多开车要一个半小时。} From Waterloo to Toronto takes an
        hour and a half by car.
  \item \zh{我从九点到十二点有课。} I have class from nine to twelve.
\end{itemize}

\subsection{\zh{先\ldots{}再\ldots} for routes}

\zh{你先坐绿线，坐三站，然后换蓝线，再坐两站就到了。} Take the green line three stops,
then change to the blue line; two more stops and you're there.

\subsection{\zh{还是} vs \zh{或者} (review)}

\begin{itemize}
  \item \zh{你坐地铁去还是坐出租车去？} Are you taking the subway or a taxi?
  \item \zh{坐地铁或者坐公交车都可以。} Either subway or bus is fine.
\end{itemize}

\subsection{\zh{每\ldots{}都\ldots}}

\zh{每} + MW + noun is almost always followed by \zh{都} before the verb.
\begin{itemize}
  \item \zh{他每天都开车上班。} He drives to work every day.
  \item \zh{每个车站都有地图。} Every station has a map.
\end{itemize}

\subsection{Conditionals: \zh{要是\ldots(的话)，就\ldots}}

\begin{itemize}
  \item \zh{要是堵车，我们就坐地铁吧。} If there's a traffic jam, let's take the subway.
  \item \zh{要是你明天有空的话，我就开车送你去机场。} If you're free tomorrow, I'll drive
        you to the airport.
\end{itemize}

\begin{important}[\zh{就} goes after the subject]
In the main clause, \zh{就} is an adverb, so it follows the subject:
\zh{我就开车送你}, not *\zh{就我开车送你}. \zh{如果} rúguǒ is a more formal synonym of
\zh{要是}.
\end{important}

\subsection{\zh{一\ldots{}就\ldots} ``as soon as''}

\begin{itemize}
  \item \zh{我一下飞机就给你打电话。} I'll call you as soon as I get off the plane.
  \item \zh{他一上车就睡觉了。} He fell asleep as soon as he got on.
\end{itemize}

\section{A subway map}

The class practised giving directions with a simplified subway map (stations
invented).

\begin{center}
\begin{tikzpicture}[st/.style={circle, draw, fill=white, thick, inner sep=2pt}, tr/.style={circle, draw, very thick, fill=white, inner sep=3pt}]
  \draw[line width=4pt, green!55!black] (0,0) -- (6,0);
  \draw[line width=4pt, blue!70] (6,2.4) -- (6,-2.4);
  \draw[line width=4pt, red!70!black] (2,-2.4) -- (2,0) -- (2,1.6) -- (4,2.4) -- (6,2.4) -- (9,2.4);
  \foreach \x/\n in {0/\zh{大学站}, 4/\zh{公园站}} {\node[st] at (\x,0) {}; \node[below, font=\small] at (\x,-0.2) {\n};}
  \node[tr] at (2,0) {}; \node[below left, font=\small] at (2,-0.1) {\zh{东门站}};
  \node[tr] (T) at (6,0) {}; \node[right, font=\small] at (6.2,0) {\zh{中心站}};
  \node[st] at (6,-2.4) {}; \node[right, font=\small] at (6.2,-2.4) {\zh{机场}};
  \node[st] at (6,-1.2) {}; \node[right, font=\small] at (6.2,-1.2) {\zh{南湖站}};
  \node[tr] at (6,2.4) {}; \node[above, font=\small] at (6,2.6) {\zh{北门站}};
  \node[st] at (9,2.4) {}; \node[above, font=\small] at (9,2.6) {\zh{火车站}};
  \node[st] at (2,-2.4) {}; \node[left, font=\small] at (1.8,-2.4) {\zh{医院站}};
  \node[font=\scriptsize, green!40!black] at (1,0.35) {\zh{绿线}};
  \node[font=\scriptsize, blue!70] at (6.6,1.2) {\zh{蓝线}};
  \node[font=\scriptsize, red!70!black] at (3,2.35) {\zh{红线}};
\end{tikzpicture}
\end{center}

\begin{example}[Directions on the map]
\zh{从大学站到机场怎么走？}---\zh{你先坐绿线，坐三站到中心站，在中心站换蓝线，再坐两站就到机场了。}
\end{example}

\section{Culture}

\begin{aside}
China's high-speed rail network (\zh{高铁}) grew from nothing in 2008 to over 30\,000
km by the end of the decade; Beijing to Shanghai (about 1300 km) takes four and a half
hours. During the Spring Festival travel rush (\zh{春运} chūnyùn) hundreds of millions
of trips are made in about forty days, and tickets sell out within minutes. Closer to
home, Waterloo's ION light rail (\zh{轻轨} qīngguǐ) opened in 2019.
\end{aside}

\section*{Review}
\begin{itemize}
  \item \zh{坐} as passenger, \zh{开} to drive, \zh{骑} for bikes.
  \item Order: subject -- time -- means (\zh{坐地铁}) -- \zh{去} + place -- purpose.
  \item \zh{从} A \zh{到} B; \zh{先\ldots{}再\ldots}; \zh{每\ldots{}都}.
  \item \zh{要是\ldots(的话)，(subject) 就\ldots}; \zh{一\ldots{}就\ldots}.
\end{itemize}
